All statements concern a frozen score vector. They do not estimate future model generations or a population of future questions. Proofs below are self-contained; the confidence-sequence construction is an application of established results, not a novelty claim.

\subsection{Decision problem}\label{theory:1}

Let \(d=(d_1,\ldots,d_N)\) belong to \([-1,1]^N\) and let \(\delta=N^{-1}\sum_i d_i\). Fix \(\epsilon\) in \([0,1)\) before examining target scores. The mutually exclusive labels are A if \(\delta>\epsilon\), E if \(-\epsilon\le\delta\le\epsilon\), and B if \(\delta<-\epsilon\). Evaluating an item reveals both model scores and therefore its difference. Known positive costs \(c_i\) refer to that pair-item evaluation. An algorithm may adaptively choose unseen items using earlier observations, metadata, and independent randomization. It must terminate by \(N\) and return one label.

Uniform \(\alpha\) correctness means that, for every allowed frozen vector \(d\), the probability of a wrong final label is at most \(\alpha\). Randomness is that of the evaluation design. A budget-limited unresolved result must be reported as unresolved, rather than E. Defining A by \(\delta>0\) instead would overlap with E; this document uses practical superiority to obtain a partition.

\subsection{Existing confidence sequences yield valid three-way decisions}\label{theory:2}

\begin{proposition}\label{app-proposition1} Suppose \(C_t\) is a \((1-\alpha)\) confidence sequence for \(\delta\): \(\mathbb P_d(\delta\in C_t\text{ for every }t)\ge1-\alpha\). Stop the first time the nonempty \(C_t\) is contained in one label region, or evaluate the full benchmark and return its exact label. Then the final label is uniformly \(\alpha\)-correct.
\end{proposition}

\begin{proof} On the simultaneous coverage event, the true \(\delta\) belongs to the region containing \(C_t\), so every certified label is correct. At census, the vector and its label are known exactly. Hence every wrong decision is contained in the CS failure event. No separate three-way Bonferroni correction is necessary. Empty sets do not yield a certificate. \end{proof}

The same event permits multiple predeclared margins or inspecting all margins after sampling, as statements about this same fixed \(\delta\). It does not license selecting a new model pair without additional multiplicity control.

\subsection{Without-replacement empirical-Bernstein construction}\label{theory:3}

Set \(x_i=(d_i+1)/2\), \(\mu=(\delta+1)/2\). Observe a uniformly random permutation \(X_1,\ldots,X_N\). Write \(S_{t-1}=\sum_{j<t}X_j\), \(R_t=N-t+1\), and \(m_t=(N\mu-S_{t-1})/R_t\). Conditional on the revealed prefix, \(\mathbb E[X_t\mid\mathcal F_{t-1}]=m_t\). Let \(p_t\) in \([0,1]\) be any predictable prediction, and choose predictable \(\lambda_t\) in \((0,1)\). Define \(\psi(\lambda)=-\log(1-\lambda)-\lambda\).

\begin{lemma}\label{app-lemma2} For each sign \(s\in\{-1,+1\}\),
\begin{equation}\label{eq:proofs-1}
M_t^s=\exp\!\left(\sum_{j\le t}\left[s\lambda_j(X_j-m_j)-\psi(\lambda_j)(X_j-p_j)^2\right]\right)
\end{equation}
is a nonnegative supermartingale starting at one.
\end{lemma}

\begin{proof} For \(y\) in \([-1,1]\) and \(0<\lambda<1\),

\begin{equation}\label{eq:proofs-2}
\exp\{\lambda y-\psi(\lambda)y^2\}\le 1+\lambda y.
\end{equation}

To verify this, expand \(\lambda y-\log(1+\lambda y)=\sum_{k\ge2}(-1)^k(\lambda y)^k/k\). For \(y\ge 0\) this is bounded by \(y^2\sum_{k\ge2}\lambda^k/k\); for \(y<0\) all terms are positive and \(|y|^k\le y^2\). The series is absolutely convergent. Apply the inequality with \(y=s(X_t-p_t)\), take conditional expectations, and multiply by \(\exp\{s\lambda_t(p_t-m_t)\}\). The result is \(\exp(-u)(1+u)\le1\), where \(u=s\lambda_t(m_t-p_t)>-1\). Multiplication over time proves the claim. \end{proof}

Let \(A_t=\sum_{j\le t}\lambda_jN/R_j\), \(Z_t=\sum_{j\le t}\lambda_j(X_j+S_{j-1}/R_j)\), and \(V_t=\sum_{j\le t}\psi(\lambda_j)(X_j-p_j)^2\). Ville's inequality for both signs gives the CS

\begin{equation}\label{eq:proofs-3}
\left[\frac{Z_t-\log(2/\alpha)-V_t}{A_t},\frac{Z_t+\log(2/\alpha)+V_t}{A_t}\right].
\end{equation}

Transform by \(2\mu-1\) for \(\delta\). This is the Waudby-Smith--Ramdas empirical-Bernstein WoR form. Our grid implementation intersects eight constant-\(\lambda\) CSs with \(\alpha/8\) each; the predictable version uses a past-only variance proxy. Hoeffding replaces the variance penalty by \(\sum_j\lambda_j^2/8\). These are alternative valid constructions, not optimized universal choices.

Intersecting intervals over earlier times preserves coverage. So does intersecting with the deterministic completion interval \([S_t/N,(S_t+N-t)/N]\). At full observation the exact mean overrides any earlier empty intersection. The numerical implementation pads bounds outward by 1e-12; the mathematical guarantee is for real arithmetic, not a formally verified floating-point implementation.

\subsection{Adaptive stratification}\label{theory:4}

Partition items into \(H\) fixed nonempty strata with weights \(w_h=N_h/N\). Construct an \(\alpha_h\) CS \([L_{h,n},U_{h,n}]\) from a uniform random permutation within each stratum, with \(\sum_h\alpha_h\le\alpha\). Adaptively select a stratum using past revealed data, and reveal its next item. Then

\begin{equation}\label{eq:proofs-4}
\left[\sum_h w_h L_{h,n_h(t)},\sum_h w_h U_{h,n_h(t)}\right]
\end{equation}

is a \((1-\alpha)\) CS for \(\delta\), using \([-1,1]\) at a zero count.

\begin{proof} The union bound gives simultaneous coverage at every local sample count in every stratum. On that event, weighting the covered stratum means gives coverage at every global time, whatever predictable interleaving was used. Thus Proposition~\ref{app-proposition1} applies. \end{proof}

This is the standard combined-stratum-CS approach. It can be conservative; it does not establish the optimality of the empirical-Neyman allocation or a general advantage of stratification. Stronger union-intersection tests already exist.

\subsection{A general indistinguishability lower bound}\label{theory:5}

\begin{theorem}[inclusion constraints]\label{app-theorem3} Fix \(\alpha<1/2\) and an algorithm uniformly \(\alpha\)-correct on a class of frozen vectors. Suppose \(d\) and \(d^{\prime}\) in this class have different labels and differ only on a set \(J\). Let \(Q\) denote the set queried before termination under \(d\). Then
\begin{equation}\label{eq:proofs-5}
\mathbb P_d(Q\cap J\ne\varnothing)\ge1-2\alpha.
\end{equation}
\end{theorem}

\begin{proof} Couple the algorithm's randomization on the two vectors, holding metadata and any side information identical. Until an item in \(J\) is queried, the two observation histories, actions, and stopping decisions coincide. The event of terminating without querying \(J\) consequently has the same coupled path on both instances. On every such path the common output is wrong on at least one of \(d\),\(d^{\prime}\), because their correct labels differ. Thus \(\mathbb P_d(Q\cap J=\varnothing)\le\mathbb P_d(\mathrm{error})+\mathbb P_{d^{\prime}}(\mathrm{error})\le2\alpha\). \end{proof}

This argument allows arbitrary adaptive, unequal-probability selection. It requires that the two constructed vectors remain admissible given side information; it cannot be applied if external information has already revealed their difference.

\begin{corollary}[cost at a boundary]\label{app-corollary3-1} Suppose \(\delta=\epsilon\). Let \(I=\{i:d_i<1\}\), and assume all vectors obtained by increasing one such coordinate to one are allowed. Then
\begin{equation}\label{eq:proofs-6}
\mathbb E_d\!\left[\sum_{i\in Q}c_i\right]\ge(1-2\alpha)\sum_{i\in I}c_i.
\end{equation}
\end{corollary}

\begin{proof} Each single-coordinate change strictly increases the mean beyond \(\epsilon\) and changes E to A. Apply Theorem~\ref{app-theorem3} to each singleton and sum \(c_i\mathbb P_d(i\in Q)\). \end{proof}

For a constant vector \(d_i=\epsilon<1\) and unit costs the bound is \((1-2\alpha)N\). At \(\alpha=.05\), any uniformly valid method must evaluate at least 90\% in expectation on that instance. For ternary paired binary scores at \(\delta=\epsilon\) and disagreement \(q\), the same bound is \((1-2\alpha)N[1-(q+\epsilon)/2]\). The negative boundary is symmetric, replacing \(d_i<1\) by \(d_i>-1\).

More generally, if \(\delta<\epsilon\), singleton constraints apply whenever \((1-d_i)/N>\epsilon-\delta\). Thus near-boundary obstructions depend on the attainable contribution of one unseen item, not solely on variance.

\subsection{Homogeneous populations away from a boundary}\label{theory:6}

\begin{theorem}[hidden exceptional items]\label{app-theorem4} Let \(d_i=b\) for every i, where \(-\epsilon\le b\le\epsilon\) and \(b<1\). Choose an integer \(1\le k\le N\) satisfying \(b+k(1-b)/N>\epsilon\), and suppose every vector formed by setting any \(k\) entries to one is admissible. Then
\begin{equation}\label{eq:proofs-7}
\mathbb E_b[\tau]\ge(N-k+1)\bigl[1-(2\alpha)^{1/k}\bigr].
\end{equation}
\end{theorem}

\begin{proof} Apply Theorem~\ref{app-theorem3} to each \(k\)-subset \(J\) and average uniformly over all such subsets, independently of the algorithm. Conditional on the realized queried set of size \(\tau\), the fraction of \(k\)-subsets avoided is \(f(\tau)=\binom{N-\tau}{k}/\binom Nk\), with zero if \(N-\tau<k\). Therefore \(\mathbb E_b f(\tau)\le2\alpha\). For integer \(t\),

\begin{equation}\label{eq:proofs-8}
f(t)\ge\max\left(1-\frac{t}{N-k+1},0\right)^k,
\end{equation}

because each factor \((N-t-j)/(N-j)\) for \(j=0,\ldots,k-1\) is at least \(1-t/(N-k+1)\) when the latter is nonnegative. The right-hand function is convex in real \(t\) for integer \(k\ge 1\). Jensen's inequality then yields \(\max(1-\mathbb E\tau/(N-k+1),0)^k\le2\alpha\), which rearranges to the result. \end{proof}

At \(b=\epsilon\) choose \(k=1\) to recover 90\% at \(\alpha=.05\). At \(b=0\), fixed positive \(\epsilon\) and large \(N\), \(k\) is about \(N\epsilon\) and the lower bound is of order \(\log(1/\alpha)/\epsilon\). This illustrates why zero observed variance does not justify immediate equivalence certification. These are elementary coupling consequences; no novelty priority is asserted for their finite-benchmark formulation. They are not asserted to be minimax-sharp away from the boundary.

\subsection{Distance-dependent consequences and limits}\label{theory:7}

\subsubsection{Distance-dependent extension: nonconstant populations on either side}\label{distance-dependent-extension-nonconstant-populations-on-either-side}

Let \(g=\bigl||\delta|-\epsilon\bigr|\). Choose a nearest decision boundary and the direction toward a different label: increase scores toward +1 for \(\delta<-\epsilon\) or \(0\le\delta\le\epsilon\), and decrease scores toward -1 for \(\delta>\epsilon\) or \(-\epsilon\le\delta<0\). The available change of coordinate i is \(a_i=1-d_i\) in the increasing direction and \(a_i=1+d_i\) in the decreasing direction. Fix \(h>0\) and let \(I_h=\{i:a_i\ge h\}\), \(M=|I_h|\). Put \(k=\lfloor Ng/h\rfloor+1\). Suppose \(k\le M\) and replacing any \(k\) coordinates of \(I_h\) by their directional endpoints is admissible given the side information. Then every uniformly \(\alpha\)-correct procedure satisfies

\begin{equation}\label{eq:proofs-9}
\mathbb E_d[|Q\cap I_h|]\ge(M-k+1)\bigl[1-(2\alpha)^{1/k}\bigr],
\end{equation}

and therefore \(\mathbb E_d[\tau]\) is at least the same quantity. With unequal costs, multiply this bound by \(\min_{i\in I_h}c_i\), and maximize over admissible \(h\). Singleton constraints can give the stronger sum-of-costs form when \(k=1\).

\begin{proof}[Proof of Theorem~\ref{main-theorem2}] Every such \(k\)-subset permits a mean change of at least \(kh/N>g\) and hence changes the label. (Starting in an outer region, overshooting to the opposite outer region still changes the original label.) Apply the inclusion constraint to every \(k\)-subset and average. Conditional on \(Q\), the fraction of subsets avoided is \(\binom{M-|Q\cap I_h|}{k}/\binom Mk\). The convex lower bound and Jensen argument of Section~\ref{theory:6} apply with \(M\) in place of \(N\). \end{proof}

The strict \(kh>Ng\) condition is sufficient on both sides and avoids endpoint-convention ambiguity; it can lose one edit when landing exactly on a closed equivalence endpoint suffices. At \(g<h/N\), \(k=1\) gives \(\mathbb E\tau\ge(1-2\alpha)M\). This describes a finite-resolution obstruction for heterogeneous vectors and both sides of the boundary. As \(g\) decreases with the capacity set held fixed, this lower bound increases stepwise. It does NOT prove that every pair of populations with ordered gaps has ordered stopping costs. The full vector determines the capacity set and the information in observations. The earlier homogeneous bound is the special case \(M=N\), \(h=1-b\) in the increasing direction. This extension is an elementary corollary of the coupling argument, not a claimed new general lower-bound technique.

\subsubsection{A complementary sufficient condition for early certification}\label{a-complementary-sufficient-condition-for-early-certification}

For uniform WoR sampling of \(d_i\) in \([-1,1]\), fix \(\alpha\) and set \(L=\log(2N/\alpha)\). A Hoeffding--Serfling interval at every \(t<N\) has radius

\begin{equation}\label{eq:proofs-10}
r_t=\sqrt{\frac{2[1-(t-1)/N]L}{t}}
\end{equation}

around the sample mean. Each interval has failure probability at most \(\alpha/N\); a union bound plus exact census yields simultaneous coverage at least \(1-\alpha\). On that event its endpoints are within \(2r_t\) of \(\delta\). For \(g>0\), the interval therefore certifies the correct region whenever \(2r_t<g\). Define

\begin{equation}\label{eq:proofs-11}
t_0=\min\left\{N,\left\lfloor\frac{8L(N+1)}{Ng^2+8L}\right\rfloor+1\right\}.
\end{equation}

The corresponding decision procedure obeys \(\mathbb P(\tau\le t_0)\ge1-\alpha\) and \(\mathbb E\tau\le t_0+\alpha(N-t_0)\). If \(t_0=N\) the statements follow from exact census. Otherwise they follow by solving \(8[1-(t-1)/N]L/t<g^2\) and applying the coverage event. This is a deliberately conservative existence bound built from an established concentration inequality, not a rate-optimal upper bound for direct betting.

Thus \(Ng^2\) much larger than \(\log(N/\alpha)\) is a sufficient regime for a small high-probability evaluation fraction, while capacity-rich populations with \(Ng<h\) have an unavoidable large fraction. These scales leave a substantial gap. A variance-based finite-population approximation suggests a crossover in \(Ng^2/\mathrm{variance}\), but neither the lower nor upper result proves a universal sharp phase transition or a variance-sensitive minimax law. The distinction between a rigorous obstruction, a sufficient condition, and an empirical crossover is essential.

No optimality theorem for the proposed bet grid, no matching upper bound for every vector, no distribution-free savings for adaptive strata, and no population-generalization guarantee are claimed. Empirical absence of violations is not a proof. Comparisons of algorithms must separate formal validity, finite Monte Carlo error, and measured cost.

\subsection{Direct directional betting and prediction-assisted signals}\label{theory:8}

For a boundary \(\mu_0\) on the transformed \([0,1]\) scale, let \(m_t(\mu_0)=(N\mu_0-S_{t-1})/R_t\). On the null \(\mu\le\mu_0\), factors \(1+\lambda(X_t-m_t(\mu_0))\) are nonnegative and have conditional expectation at most one when \(0\le m_t(\mu_0)\le1\) and \(\lambda\) belongs to the declared (0,1) grid. Use unit factors when that candidate conditional mean is outside \([0,1]\). If the bounded completion interval excludes the entire null, rejection is deterministically sound. The upper-direction factors use the minus sign and null \(\mu\ge\mu_0\). A fixed average of the eight products is a supermartingale under each respective null. Its first crossing of \(2/\alpha\) is a level-\(\alpha/2\) test.

Use four such tests: lower and upper directions at both -\(\epsilon\) and +\(\epsilon\) (transformed to \([0,1]\)). A requires rejection of the lower-direction null at +\(\epsilon\); B requires rejection of the upper-direction null at -\(\epsilon\); E requires both inner-direction rejections. Conflicting evidence yields no certificate. If E is true, an error requires either of the two outer rejections. If A is true, an error requires an upper-direction rejection at +\(\epsilon\) (for false E) or at -\(\epsilon\) (for false B). If B is true the two lower-direction rejections play the corresponding role. In every case a union bound over two level-\(\alpha/2\) tests suffices. Strict tests may not certify a boundary before census; census returns the closed-region label exactly.

For the prediction-assisted adaptation, the charged pilot reveals a set O. Subsequent selection uses predictable \(q_i>0\) on the remaining set \(J\), with \(\sum_i q_i=1\), and predictable predictions \(p_i\) in \([0,1]\). Define the CELEUS finite-pool signal

\begin{equation}\label{eq:proofs-12}
Y=\frac{\sum_{i\in O}x_i+\sum_{i\in J}p_i}{N}+\frac{x_I-p_I}{Nq_I}.
\end{equation}

Conditional on the past, \(\mathbb E[Y]=\mu\), because the \(q\)-weighted correction averages to the remaining total residual divided by \(N\). Put \(\mathrm{base}=[\sum_{i\in O}x_i+\sum_{i\in J}p_i]/N\). Valid predictable bounds are \(a=\mathrm{base}-\max_{i\in J}p_i/(Nq_i)\) and \(b=\mathrm{base}+\max_{i\in J}(1-p_i)/(Nq_i)\). Hence \(Z=(Y-a)/(b-a)\) is in \([0,1]\) and has conditional expectation \((\mu-a)/(b-a)\). Apply the preceding directional factors to Z and candidate \((\mu_0-a)/(b-a)\). Unit factors handle candidates outside the signal support. The same two-event argument proves \(\alpha\) correctness.

The estimator signal is adopted from CELEUS, not proposed as new. The implemented predictor is a nearest historical response profile selected using 64 charged pilot pair-items; reference selection and sampling weights are frozen thereafter. It is an explicitly documented adaptation, not a reproduction of the paper's surrogate-model experiments. Positive fixed weights generate a size-biased WoR permutation through independent exponential clocks; at each reveal the conditional selection probability is the item's weight divided by the remaining sum. Nonuniform selection is therefore accounted for exactly. No assumption that the predictor is accurate appears in the validity proof.

\subsection{\texorpdfstring{Joint stratified tests without \(\alpha/H\) splitting}{Joint stratified tests without \textbackslash alpha/H splitting}}\label{theory:9}

This is a specialization of established union-intersection stratified inference, not a claim to invent that approach. Let \(\mu_h\) be the fixed transformed mean in stratum \(h\), and \(w_h=N_h/N\). For the first \(n_h\) observations of that stratum use the unweighted local statistics \(A_h=\sum_j N_h/(N_h-j+1)\), \(Z_h=\sum_j[X_{h,j}+S_{h,j-1}/(N_h-j+1)]\), and \(V_h=\sum_j(X_{h,j}-p_{h,j})^2\). Here \(p_{h,j}\) is a predictable local forecast in \([0,1]\). These remove the \(\lambda\) and \(\psi\) weights from the statistics in Section~\ref{theory:3} because the weights are applied explicitly below. Set all three statistics to zero before the first observation. For a fixed \(\lambda\) in \((0,1)\) and direction \(s\in\{-1,+1\}\), the process

\begin{equation}\label{eq:proofs-13}
M_t(\lambda,\mu,s)=\exp\!\left\{s\lambda\sum_h[Z_h-A_h\mu_h]-\psi(\lambda)\sum_hV_h\right\}
\end{equation}

is a nonnegative supermartingale at the true vector of stratum means under any predictable allocation that samples uniformly within the selected remaining stratum. To verify this, on a step in \(h\) write m for its true remaining mean. The increment of \(Z_h-A_h\mu_h\) is \(X-m\). The scalar inequality proved in Section~\ref{theory:3} gives

\begin{equation}\label{eq:proofs-14}
\begin{aligned}\mathbb E[\exp\{s\lambda(X-m)-\psi(\lambda)(X-p)^2\}\mid\mathcal F_{t-1}]\\\le \exp\{s\lambda(p-m)\}[1+s\lambda(m-p)]\le1.\end{aligned}
\end{equation}

Here \(|X-p|\le1\), the scalar inequality applies to \(s(X-p)\), and \(1+u\le\exp(u)\) proves the second inequality with \(u=s\lambda(m-p)\). All other strata's factors stay fixed. Predictability of the selected stratum establishes the claim in the global filtration. An average over the fixed \(\lambda\) grid is also a supermartingale.

At each inspection let \(B_t\) be the Cartesian product of the deterministic completion intervals \([S_h/N_h,(S_h+N_h-n_h)/N_h]\). For the lower-direction composite null \(w\cdot\mu\le m_0\), minimize the mixture over \(B_t\) intersected with that null. If this set is empty, rejection is deterministically valid. Otherwise the minimizing vector maximizes \(A\cdot\mu\) over that set: every \(\lambda\) is positive, so the same vector minimizes every component of the mixture. Since \(A_h\ge 0\), a maximizer can be taken on \(w\cdot\mu=\min(m_0,w\cdot\mathrm{hi})\). Starting at lo, fill coordinates in decreasing order of \(A_h/w_h\) until that weighted sum is attained. This fractional-knapsack solution is exact: transferring weighted mass from a smaller ratio to a larger unsaturated ratio cannot decrease the objective, which proves optimality by successive exchanges. The upper-direction null \(w\cdot\mu\ge m_0\) instead minimizes \(A\cdot\mu\), fills in increasing ratio order, and uses boundary \(\max(m_0,w\cdot\mathrm{lo})\).

Reject when this infimum reaches \(2/\alpha\) at any inspection. The infimum process itself need not be a supermartingale. Nevertheless, whenever the null is true its feasible set contains the true vector on every path. Crossing of the infimum therefore implies crossing of the true-vector mixture; Ville's inequality bounds its probability by \(\alpha/2\). The four directional tests at the two practical margins give \(\alpha\)-correct ternary decisions by the same two-event argument as Section~\ref{theory:8}. Inspections every 16 observations restrict the rejection times and preserve validity; all observations between inspections are charged. Exact census supplies the terminal label. Adaptive allocation does not invalidate the argument, but no efficiency optimality follows.

Implementation checks compare 300 optimizations (150 random boxes, both directions) against an independent linear-programming solver and enumerate all 729 ordered ternary populations split into two strata, under each of two allocations. These checks supplement the proof; small-population tests alone cannot establish general error control.


